\section{HiGTree: a camada de malha}

Trinta e tr\^es fontes, 15\,869 linhas. \'E a biblioteca sobre a qual tudo o mais
se ap\'oia, e ela n\~ao sabe nada de Navier--Stokes.

\begin{adjustbox}{max width=\linewidth}
\pgfplotstabletypeset[
  string type,
  columns/modulo/.style={column name={m\'odulo}},
  columns/linhas/.style={column name={linhas}},
  columns/papel/.style={column name={papel}},
  every head row/.style={before row=\toprule, after row=\midrule},
  every last row/.style={after row=\bottomrule},
]{\dat{higtree-modulos.dat}}
\end{adjustbox}

\paragraph{A divis\~ao que mais importa \'e \texttt{domain} contra
\texttt{pdomain}.} O primeiro \'e o dom\'inio serial --- c\'elulas, facetas,
estênceis, mapeadores. O segundo \'e o distribu\'ido, e \'e onde vivem a
\emph{franja} (as c\'opias de c\'elulas de outro \emph{rank}), a
sincroniza\c{c}\~ao e os tipos derivados de MPI. Juntos s\~ao
5\,115 linhas, quase um ter\c{c}o da biblioteca: a distribui\c{c}\~ao \'e o
assunto mais caro dela.

\paragraph{Solvers lineares: cinco implementa\c{c}\~oes atr\'as de uma
interface.} PETSc, HYPRE, SOR, centralizado e um de depura\c{c}\~ao que
simplesmente grava o sistema. O de PETSc \'e o usado na pr\'atica. Existe
tamb\'em um \texttt{solver-\allowbreak{}petsc.\allowbreak{}camila.\allowbreak{}c}, variante paralela do principal ---
duas c\'opias de 576 e 555 linhas que divergir\~ao na primeira mudan\c{c}a de
protocolo.

\paragraph{\texttt{wls.c} \'e o cora\c{c}\~ao do m\'etodo.} Sem os m\'inimos
quadrados ponderados, a \'arvore n\~ao graduada n\~ao fecha: \'e ele que produz
valor onde o estêncil regular n\~ao existe. As 624 linhas s\~ao pequenas para o
peso que carregam.
